% Section 9 (group D): open problems and testable statements.
% Sources: Companion rounds v4.10.0-r05, r08, r18, r21, r28, r30, r31, r55, r56, r59, r60, r72--r75;
% wiki/Vorhersagen-und-Tests.md, wiki/Offene-Probleme.md, wiki/Gesamtbild.md.
\section{Open problems and testable statements}
\label{sec:open}

\subsection{The open problems}
\label{sec:open-problems}

The Companion names five open problems OP1--OP5 and two closure questions, C5 and C6/T2. None of them is closed
\cite{QSC}. Their present formulation is best read through one observation of \rnd{r30} and \rnd{r55}: every number that
OP2--OP5 supply today is a function of the single number of OP1, $\bsym=0.5494$, but it enters each problem through a
separately set identification, four of them chosen after the data were known (\rnd{r30}, remark r30.6). Each problem
also keeps a remainder that does not depend on the number (\rnd{r55}, remark r55.4).

\paragraph{OP1: the origin of $\bsym$.} \Open\ The frame offset is computed from a counting rule that is postulated, not derived (\rnd{r67}):
\begin{equation}
\bsym=\sqrt{\frac{\gamma_L-1}{\gamma_L+1}}=0.5494,\qquad \gamma_L=\frac{\gamma_M}{\gamma_T}=1.8647,
\label{eq:open-beta}
\end{equation}
where $\gamma_M=0.237533$ is the
Barbero--Immirzi value of the counting rule with all spins and weight two (assigned to frame M) and
$\gamma_T=\ln2/(\pi\sqrt3)=0.127384$ that of the rule with $j=\tfrac12$ only (assigned to frame T). The assignment is
equivalent to a frame-dependent area quantum (\rnd{r10}). Unitary boosts leave the area spectrum fixed, the Holst duality
rotation does not map $\gamma_M$ to the needed value, and the non-unitary frame map $J_B$ is a conjugation that leaves every
spectrum fixed (\rnd{r30}, \rnd{r31}, \rnd{r52}). OP1 therefore rests on the counting rules alone: why M counts every spin with
weight two, the extremal doublet of each puncture (the premise called U2) and why T counts only $j=\tfrac12$. Both premises and the Casimir area $\sqrt{j(j+1)}$ form the counting postulate adopted by the author on 26 November 2026 (\rnd{r67});
with the area $j+\tfrac12$ the same rule would give $1/\sqrt3$. The postulate replaces two literature values by one rule with three named
parts and leaves OP1 open in the form ``derive the counting postulate''. A solution would supply the
number but none of the five identifications (\rnd{r55}, r55.1). Section~\ref{sec:beta} gives the details.

\paragraph{OP2: the Dirac operator on the antisymmetric sector.} \Open\ Existence and self-adjointness of the operator $D_-$
on the antisymmetric sector were derived from an earlier manuscript of the author, which has not been a source of proof since v4.8.0-r04.
Its own remainder is that proof; its identification with the frame offset (spectral offset $=\bsym$) is one of the five
set translations (\rnd{r30}, r30.6).

\paragraph{OP3: generations and the residual bend.} \Open\ The Companion does not derive three generations, has no ban on
a fourth, and does not derive Koide's $K=2/3$; the overlap $\tfrac12$ is an input (Section~\ref{sec:B-koide}). Its number-dependent part is
the residual bend $\Phi_e=\pi/4-\gamma_M/2$ of the lepton phase, tested by H102. The orientation of the octant step is a
time orientation of the readout that no candidate of the Companion fixes (H118); on this the author decided on 22~November
2026 to let OP3 rest until a new idea appears (\rnd{r31}, r31.2). The Albert-algebra route gives a family $SO(3)$ broken by charge
exactly to $A_4$, with $V_{\mathrm{CKM}}=\mathbf 1$ for exact $A_4$, but no numbers (\rnd{r60}; Section~\ref{sec:B-particles}).
After \rnd{r60} two of these questions were answered without closing OP3 (tests H138--H156 of 25 and 26~November 2026, partly
not blind; recorded in Companion rounds \rnd{r62}, \rnd{r63}). The three $G_2$-stable octonion copies in the sedenions share one colour
$su(3)$ and have no common charge: they are three readings of one generation, not three generations (H138, H140); the number of generations stays an input. In an
amplitude readout $K=1/(3\cos^2\theta)$, so $K=\tfrac23$ holds exactly at the mixing angle $\theta=\pi/4$, the fixed point of
a swap read as the exchange of matter and tachyon frame (H148, H149); the swap symmetry is a postulate that the algebra
does not supply (H156), and the phase $\delta=2/9$ remains a fit. Rounds \rnd{r73} and \rnd{r74} reduce it to one set sign and one
rotation angle, $\varphi_\ell=0.0395773$, which no route tried derives and which stays a measured input (Section~\ref{sec:B-lepton-phase}).

\paragraph{OP4: the origin of the mixing angles.} \Open\ Several constructions reproduce the mixing data, but as fits.
The translations of the number are $|U_{e2}|^2=\bsym^2$ (H9) and the mask parameter $\varepsilon=\eta$ (H46). The own
remainder is the $\mathbb Z_2$ of the neutrinos (TM2), $\delta_{\mathrm{CP}}$, the CKM matrix and the assignment of gauge
charges (\rnd{r55}). One relative sign follows from the structure once the same side of the multiplication is taken
for the family in both sectors (\rnd{r65}): the relative orientation of the family $su(2)$ on leptons and quarks (H125 stage A, not blind), which under a
further setting S1 becomes the prediction $\sin\delta_{\mathrm{PMNS}}<0$ (\rnd{r59}).

\paragraph{OP5: the cosmological constant and the vacuum balance.} \Open\ By an author decision OP5 has three parts
(\rnd{r05}, r05.3): the core closure of the source--return route, locally closed under the guards and not claimed
globally; the $\Lambda$ readout $\Omega_M=\bsym^2$ (H96, open); and the action programme W1/W2, frozen
(Section~\ref{sec:act}). The own remainder, the absolute value of $\Lambda$, is an input for the named class of
constructions of the Companion: no equation in it selects $\Lambda$ and the class has no scale of its own; the readout
$\Omega_M=\bsym^2$ only trades $\Lambda$ for the epoch (\rnd{r56}). This is a negative result for the Companion as it
stands, not for every extension (Section~\ref{sec:lop5}).

\paragraph{C5.} \Open\ In its old form C5 asked, in the measure of one frame, about a passage between two frames. It was
restated in \rnd{r21} as three questions. The transfer question $C5_{\mathrm{tr}}$, typed by frame, is, within the reading of \rndold{r10}, a relocation: an
empty compensator port is a full partner port, since $C5_M+C5_T=1$; nothing vanishes. The frame-free photon-point
question $C5_P$ has a positive answer as a statement about fields: free time evolution reaches $|s|=0$ at a finite time
on a codimension-one family. Only the bridge question $C5_{\mathrm{br}}$, whether reaching the photon point discharges the
obstruction, is open; under the original definition of the score its answer is negative (r21.12), and the assertion
$C5=0$ is not made.

\paragraph{C6/T2 and the guards.} The global closure of all readouts, C6/T2, is \emph{not claimed}. The Companion
asserts only local, frame-typed closures (thesis T15). Its five guards are unchanged: every round up to \rnd{r75} records
$R_{\mathrm{glob}}=(0,1,1,1,1)$, $\Gamma_{\mathrm{guard}}=(1,1,1,1,1)$ and the end bound $\beta_4=1/256$ unspent, and states that it
implies no exact finite C5 zero closure, no exact all-readout terminality, no C6/T2, no global photon-clock theorem and no
global all-readout closure. Round \rnd{r72} restates the fourth guard, unchanged in content and value, as ``no global photon-clock
theorem''.

\subsection{Testable statements}
\label{sec:open-tests}

The Companion fixes each prediction in a pre-registration note (hypothesis, data, decision criterion) committed before
the decisive data exist; a value counts as tested only with a measurement after the stated cutoff, and earlier data are
given for comparison only. Table~\ref{tab:open-tests} lists the predictions that are open, together with H2. The status is
that of 25~November 2026; the dated list of this version, with the neutrino masses of H227 and the mixing readings recorded after
that date, is Table~\ref{tab:open-predictions}. None of the open tests is decided. One further open pre-registration is in no Companion round up to
\rnd{r61} and not in the table: H105 (horizon area steps $4\ln2\,\ell_P^2$ and ringdown echoes under an added setting V1,
deadline 31.12.2028; recorded in \rnd{r64}); its value is not specific to the construction.

\begin{table}[ht]
\centering
\footnotesize
\renewcommand{\arraystretch}{1.15}
\begin{tabularx}{\textwidth}{@{}l>{\raggedright\arraybackslash}p{0.27\textwidth}>{\raggedright\arraybackslash}X>{\raggedright\arraybackslash}p{0.25\textwidth}@{}}
\toprule
& statement & decisive data; refuted if & status and remarks\\
\midrule
H2 & normal ordering, $\sum m_\nu=0.0592$\,eV (neutrinos as light-like triple) & mass ordering and $\sum m_\nu$ &
  compatible, weakly discriminating; ordering not yet determined (\rnd{r18})\\
H9 & $\sin^2\theta_{12}=\bsym^2/\cos^2\theta_{13}=0.3086$ & JUNO with $\sigma\le0.004$ &
  open; found after the fact; JUNO 2026: $0.3036\pm0.0064$ (\rnd{r31})\\
H15 & $m_W=80.3615\pm0.0018$\,GeV & first world average after 21.10.2026 &
  open; PDG 2026 ($80.3625\pm0.0077$) predates H15 and does not count (\rnd{r04})\\
H42 & $\sin^2\theta_{23}\le0.54$ (normal ordering) & first global analysis or DUNE/Hyper-K result after 24.10.2026; refuted if the
  $2\sigma$ interval lies wholly above $0.54$ & open; independent of $\gamma_L$ (\rnd{r08})\\
H46 & $\sin^2\theta_{23}\in[0.49,0.50]$ under $\varepsilon=\eta$ & same data; refuted if the $2\sigma$ interval misses $[0.49,0.50]$ &
  open\\
H96 & $\Omega_M=\bsym^2=0.30185$ today & first flat-$\Lambda$CDM $\Omega_m$ from DESI DR3 (BAO alone) or Euclid; refuted beyond $2\sigma$ &
  open; DESI DR2 $0.2975\pm0.0086$ known, does not count (\rnd{r05})\\
H100 & $R=\Delta m^2_{21}/\Delta m^2_{31}$ in the next NuFIT fit: A $0.030839$, B $0.030087$, C $0.030375$ & first NuFIT version after 6.1;
  confirmed if A is refuted and B or C compatible ($|z|\le2$); lapses without a new version by end of 2027 & open (\rnd{r12})\\
H102 & $m_\tau=1777.134$\,MeV from $\Phi_e=\pi/4-\gamma_M/2$ (exact Koide: $1776.969$) & first PDG edition after 13.11.2026, at the
  latest end of 2028; refuted if $|z|>2$ & open; $+2.3\sigma$ against PDG 2024, $1776.93\pm0.09$, unchanged in PDG 2026 (\rnd{r28}.3); exact Koide lies at $+0.4\sigma$\\
H112 & $\sin^2\theta_{12}\cos^2\theta_{13}=\Omega_m$, the same number at both places, whatever it is & data of H9 and H96 together;
  refuted if $|z|\ge3$; deadline 31.12.2029 & open; weak (refutes only from $|\Delta\beta^2|\gtrsim0.022$)\\
H125 & $\sin\delta_{\mathrm{PMNS}}<0$ (relative CP sign of leptons and quarks, under S1 and the same-side setting) & first analysis fixing the sign at $\ge3\sigma$;
  deadline 31.12.2035 & open; stage A confirmed, \emph{not blind}; prior probability $\tfrac12$ (\rnd{r59}); NuFIT fits already
  lean to $\sin\delta<0$\\
\bottomrule
\end{tabularx}
\caption{Open pre-registered predictions of the Companion rounds, and H2, status of 25 November 2026 (wiki page ``Vorhersagen und
Tests'' and the cited rounds). H15, H46 and H96 rest on one common core, the assignment of counting rules to the frames;
H102 is a joint test of the core reading, the $U(1)$ counting and the area spectrum $\sqrt{j(j+1)}$.}
\label{tab:open-tests}
\end{table}

Three qualifications belong to every entry. First, the values of H9, H15, H46, H96 and H102 depend on the number of OP1, and
the chance estimate of Section~\ref{sec:beta} applies to them: as long as the pair of counting rules is chosen rather than derived, its
hits are compatible with chance. H15, H46 and H96 are not three independent tests; together with H42 the four older
predictions form two independent groups (\rnd{r08}). Second, H9 and H102 were found after the older data were seen; they
are blind only against future measurements. Third, H112 tests the weaker statement that the same number stands at both
places.

\Open\ \textbf{No blind prediction of the Companion has so far been confirmed against measurement data.} The confirmed
tests without the mark ``not blind'' all check inner mathematics (\rnd{r05}, remark r05.2), and no later test, up to
\rnd{r75}, has changed this. The predictions above and in Section~\ref{sec:open-predictions} are the places where that can change.

\subsection{Predictions of this version (as of 28 November 2026)}
\label{sec:open-predictions}

This version of the article (1.2.0) states the model of Companion v5.6.0 (Version DOI
\href{https://doi.org/10.5281/zenodo.24632437}{10.5281/zenodo.24632437}; predecessor v5.5.0, 10.5281/zenodo.24563811), as of
28~November 2026. Table~\ref{tab:open-predictions} collects its predictions. The full dated register, with deciding experiments, deadlines
and refutation criteria, is round \rnd{r75} of the Companion, which is authoritative; the binding criteria stand in the advance notes
of the tests. The last column says what each value rests on: \emph{derived} means computed from the stated inputs, \emph{set} a postulate
or a setting, \emph{measured} a measured input or a scale fitted to data, \emph{post hoc} a relation found after the data were known.
Common to all lepton rows: the Brannen form \eqref{eq:B-brannen-sector} is set; $K=\tfrac23$ holds under the postulate of equal standing
of the frames; $\varphi_\ell=0.0395773$ is a measured input; the sign of the branch choice is a setting; the scales $\mu_e$, $\mu_\nu$
are fitted. Common to the rows that use $\bsym$: the value $\bsym=0.5494056410$ is computed from the counting postulate of \rnd{r67}
and is not derived.

\begin{table}[htp]
\centering
\footnotesize
\renewcommand{\arraystretch}{1.12}
\begin{tabularx}{\textwidth}{@{}>{\raggedright\arraybackslash}p{0.15\textwidth}>{\raggedright\arraybackslash}p{0.26\textwidth}>{\raggedright\arraybackslash}p{0.12\textwidth}>{\raggedright\arraybackslash}X@{}}
\toprule
quantity & prediction & source & rests on\\
\midrule
mass ordering & normal & H227, H213; \rnd{r73} & derived from $\varphi_\ell$ (measured); branch sign and Brannen form set\\
$\sum m_\nu$ & $59.30\pm0.24$\,meV & H227; \rnd{r73} & derived; $\varphi_\ell$ measured, $\mu_\nu$ fitted to $\Delta m^2_{31}$; $K=\tfrac23$ postulate; form and sign set\\
$m_1$, $m_2$, $m_3$ & $0.379$, $8.81$, $50.11$\,meV & H227; \rnd{r73} & as for $\sum m_\nu$\\
$m_\beta$ & $8.87$\,meV & H227; \rnd{r73} & as for $\sum m_\nu$; mixing data measured\\
$m_{\beta\beta}$ & $1.3$ to $4.0$\,meV (Majorana phases free) & H227; \rnd{r73} & as for $m_\beta$\\
$R=\Delta m^2_{21}/\Delta m^2_{31}$ & $0.030839$ (data $+1.8\sigma$; not blind) & H99, H100, H227 & derived from $\varphi_\ell$ (measured); sector assignment H99 and form set\\
$\sin^2\theta_{12}$ & three exclusive readings: $0.3086$; $0.3019$; $0.318$ to $0.3195$ & H9; H176; H206--H208 & H9: $\bsym^2$ (postulate), found post hoc; H176: one-number candidate with set readouts, selected post hoc; H206--H208: column $|U_{e1}|^2=\tfrac23$ derived, readout set\\
$\theta_{12}$--CP correlation & $\sin^2\theta_{12}\gtrsim0.315$ with $|\sin\delta|\approx1$, or $\sin^2\theta_{12}$ near $0.304$ with $|\sin\delta|\approx0.65$ & H206--H208 (JUNO with DUNE) & structural readout set; additional term allowed, not derived; correlation found post hoc (H208)\\
$\sin^2\theta_{23}$ & $0.4954$ (H176); $\le0.54$ (H42); $[0.49,0.50]$ (H46) & H176; H42, H46 & set readouts; $\bsym$ from the postulate (H46)\\
$\sin^2\theta_{13}$ & $0.0223$ (H176); $0.0216$ without parameter (H207) & H176, H207 & set readouts; H207 partly not blind\\
$\delta_{\rm CP}$ & $|\sin\delta|\approx0.977$ and $\sin\delta<0$, together about $282^\circ$ & H176; H125 & modulus as for $\theta_{23}$; sign rests on setting S1 and the same-side setting\\
$\Omega_m$ (flat $\Lambda$CDM) & $\bsym^2=0.30185$ & H96 & $\bsym$ from the postulate; identification $\Omega_M=\bsym^2$ set\\
$m_W$ & $80.3615\pm0.0018$\,GeV & H15 & $\gamma_L$ from the postulate; $\cot\theta_W=\gamma_L$ set, found post hoc; $m_Z$ measured\\
$m_\tau$ & $1776.969$\,MeV (Koide; not blind); alternative $1777.134$\,MeV & \rnd{r63}, H227; H102 & $K=\tfrac23$ postulate, $m_e$, $m_\mu$ measured; the alternative rests on the bent-octant placeholders, found post hoc\\
\bottomrule
\end{tabularx}
\caption{Predictions of version 1.2.0, as of 28 November 2026 (Companion v5.6.0, round \rnd{r75}). The rows are not independent: H15,
H46 and H96 share the counting core; H9, H96 and H112 share $\bsym^2$; the three readings of $\theta_{12}$ exclude one another.}
\label{tab:open-predictions}
\end{table}

\paragraph{Version rule.} By a standing rule of the author (28 November 2026; \rnd{r75}.1) every public version records its predictions
with date and DOI. These numbers are fixed with this version and are not readjusted after a new measurement. A correction appears only in
a later version, with its reason stated (progress of the model or an error). The version that was fixed before the measurement counts as
the test, and earlier versions stay visible under their DOIs. Several rows fall together: an established inverted ordering or a measured
$\sum m_\nu<58$\,meV refutes the mass balance and with it the Brannen form for neutrinos; JUNO separates the three readings of
$\theta_{12}$ and, with DUNE, tests the correlation; a deviation of $\Omega_m$ or $m_W$ bears on the counting postulate; the two values of
$m_\tau$ separate only at an uncertainty of about $0.05$\,MeV. The register adds no derivation, no guard is spent, and OP1--OP5, C5 and
C6/T2 remain open.
